\documentclass[11pt]{article}
\usepackage[T1]{fontenc}
\usepackage{lmodern}
\usepackage{amsmath,amssymb,amsthm}
\usepackage[margin=1.05in]{geometry}
\usepackage{microtype}
\usepackage[colorlinks=true,linkcolor=blue,citecolor=blue,urlcolor=blue]{hyperref}
\numberwithin{equation}{section}
\newtheorem{theorem}{Theorem}[section]
\newtheorem{lemma}[theorem]{Lemma}
\newtheorem{corollary}[theorem]{Corollary}
\newcommand{\CC}{\mathbb C}
\newcommand{\Aff}{\mathbb A}
\DeclareMathOperator{\grad}{\nabla}
\title{An explicit noncoordinate polynomial\\with affine three-space zero fibre}
\author{OpenAI}
\date{September 24, 2026}
\hypersetup{
  pdftitle={An explicit noncoordinate polynomial with affine three-space zero fibre},
  pdfauthor={OpenAI}}

\begin{document}
\maketitle

\begin{abstract}
We construct an explicit counterexample to the Abhyankar--Sathaye conjecture:
a noncoordinate polynomial in four complex variables whose zero fibre is affine
three-space. Adjoining variables gives counterexamples in every ambient
dimension at least four.
\end{abstract}

\section{Introduction and statement}\label{sec:introduction}

The Abhyankar--Sathaye conjecture asks whether a polynomial
$f\in\CC[z_1,\ldots,z_n]$ satisfying
\[
  \CC[z_1,\ldots,z_n]/(f)\cong\CC^{[n-1]}
\]
must be a coordinate, that is, the first component of a polynomial automorphism
of $\Aff^n_{\CC}$.
Here $\CC^{[m]}$ denotes a polynomial ring in $m$ independent variables, and
the displayed isomorphism is required to fix $\CC$.
Geometrically, the question asks whether
every hypersurface embedding $\Aff^{n-1}_{\CC}\hookrightarrow\Aff^n_{\CC}$ can
be carried to a coordinate hyperplane by an ambient polynomial automorphism.

For $n=2$, the question is answered affirmatively by the
Abhyankar--Moh--Suzuki Theorem~\cite{AM,Suzuki}.
The higher-dimensional problem belongs to a family of questions about
polynomial maps with affine-space fibres. Sathaye~\cite[pp.~1--2]{Sathaye}
distinguished the coordinate, or epimorphism, question from the weaker question
whether an affine-space zero fibre forces all fibres to be affine space.
A coordinate polynomial defines a projection after an automorphism, so its
fibres are all isomorphic to $\Aff^{n-1}_{\CC}$; an isomorphism of the zero fibre
alone supplies no such global trivialization.
Kraft~\cite[\S4]{Kraft} surveys the embedding problem and its relation to
algebraic group actions. Among the early affirmative results,
Sathaye~\cite{Sathaye} proved rectifiability for planes defined by
polynomials linear in one of three variables in characteristic zero, and
Russell~\cite[Theorem~2.3]{Russell} extended this result to arbitrary fields.

In four variables, Kaliman, V\'en\'ereau and Zaidenberg~\cite[Theorem~4.2(a)]{KVZ}
and Maubach~\cite[Corollaries~4.2--4.3]{Maubach} established affirmative results
for equations $a(x)y+b(x,z,t)$ with $a\ne0$: an affine-three-space zero fibre
can be straightened by an ambient automorphism. Here $x,y,z,t$ are independent
variables, and the coefficient of $y$ depends on only one of them.
Further algebraic treatments appear in work of El~Kahoui, Essamaoui and
Ouali~\cite{EEO} and Ghosh, Gupta and Pal~\cite{GGPthreefold}.
Ghosh, Gupta and Pal also study linear hypersurfaces with multivariable
coefficients under additional hypotheses on the coefficients and their
divisors~\cite[Theorem~4.11]{GGP}.
Ghosh~\cite[Theorem~1.4]{Ghosh} proves an affirmative embedding theorem in
characteristic zero for
$a(x_m)b(x_1,\ldots,x_{m-1})y+f(z,t)+x_m$, where the displayed variables are
independent, $b\ne0$, and $a$ has at least two distinct roots over the algebraic
closure.
These results illustrate the role of restrictions on the defining equation;
an abstract polynomial presentation of the quotient does not itself provide
an automorphism of the ambient polynomial ring.

We give a counterexample in ambient dimension four, and hence in every
ambient dimension at least four. The three-variable case is not addressed.

The construction starts from the cusp parametrization $(u^3,-u^2)$ and
perturbs its two entries by independent multiples of $h$.
In $R=\CC[h,u,v,w]$, put
\begin{equation}\label{eq:xyz}
  \begin{aligned}
    x&=u^3+hv,\qquad y=-u^2+hw,\\
    s&=S(h,u,v,w)
       :=2u^3v+3u^4w+h(v^2-3u^2w^2)+h^2w^3.
  \end{aligned}
\end{equation}
Since $x^2+y^3$ vanishes modulo $h$, it is divisible by $h$ in $R$.
The displayed polynomial $s$ is its quotient, so $x^2+y^3=hs$.
Define $p$ and $F$ by
\begin{equation}
  p(x,y,s)=-2s^2x+3sy^2-3s^3y,
  \qquad F=h-p(x,y,s)-1. \label{eq:F}
\end{equation}
Thus $F$ is a polynomial in the four independent variables $h,u,v,w$.

\begin{theorem}\label{thm:main}
The polynomial $F$ defined in \eqref{eq:xyz}--\eqref{eq:F} satisfies
\[
  R/(F)\cong\CC^{[3]},
  \qquad
  \grad F(2,0,-1/2,1/2)=0.
\]
Consequently $F$ is not a coordinate of $R$.
\end{theorem}

The critical point $P=(2,0,-1/2,1/2)$ lies on the fibre $F=-1$.
The zero fibre is smooth, since it is isomorphic to affine three-space.

\begin{corollary}\label{cor:dimension}
For every integer $n\ge4$, there is a polynomial
$f\in\CC[z_1,\ldots,z_n]$ such that
$\CC[z_1,\ldots,z_n]/(f)\cong\CC^{[n-1]}$ and $f$ is not a coordinate.
\end{corollary}

Shpilrain and Yu~\cite[\S1]{SY} used the ambient critical locus to distinguish
isomorphic hypersurfaces with inequivalent embeddings. They asked whether the
gradient must be nowhere zero when the zero fibre is isomorphic to a coordinate
hyperplane. Theorem~\ref{thm:main} gives a negative answer in ambient dimension
four, and Corollary~\ref{cor:dimension} preserves this obstruction after adjoining
variables. The distinction is between the intrinsic geometry of one fibre and
the behaviour of its defining polynomial on the ambient space.

The main algebraic ingredient is Lemma~\ref{lem:lift}. It turns a presentation
over a commutative base ring $B$ with $x^2+y^3=hs$ and $(h,y)=B$ into a polynomial
extension of that base. Two compatible linear equations, with coefficients
generating the unit ideal, give explicit inverse maps even when the base has
zero divisors. In our example, a translation identifies the auxiliary algebra
cut out by $x^2+y^3=hs$ and $h=1+p(x,y,s)$ with a polynomial ring in two
variables; the lemma supplies the third variable. This proves that the zero
fibre is affine space. An identity in the original four-variable ring then
exhibits a critical point on the fibre over $-1$, which rules out a coordinate.

The polynomials in \eqref{eq:xyz} also have an invariant-theoretic
interpretation. The triangular derivation
$D=h\partial_u-3u^2\partial_v+2u\partial_w$ of $R$ annihilates $h,x,y,s$.
They arise in the triangular monomial kernel calculation of
Maubach~\cite[proof of Theorem~3.2, case~(3)]{MaubachMonomial}.
The proof below uses the cusp relation directly through Lemma~\ref{lem:lift}.

Section~\ref{sec:lifting} proves the lifting lemma over an arbitrary commutative
base ring. Section~\ref{sec:mainproof} identifies the auxiliary algebra,
establishes the zero-fibre isomorphism, and proves the critical-point obstruction
and its persistence after adjoining variables. Appendix~\ref{app:coordinates}
records the resulting polynomial coordinates and their inverse.

\section{An elementary lifting lemma}\label{sec:lifting}

The next lemma provides the polynomial parameter used to identify the zero
fibre. Its unit-ideal hypothesis allows us to solve two compatible linear
equations without dividing by either coefficient. This point matters because
the coefficients need not be units, and the base ring may have zero divisors.

In the following statement, $S(H,U,V,W)$ denotes the universal polynomial
specified in \eqref{eq:xyz}, with the indicated arguments. All rings are
commutative with identity.

\begin{lemma}\label{lem:lift}
Let $B$ be a ring and let $h,x,y,s\in B$ satisfy
\begin{equation}
  x^2+y^3=hs,\qquad (h,y)=B. \label{eq:base}
\end{equation}
Then the $B$-algebra
\[
  C=B[U,V,W]/\bigl(U^3+hV-x,\,-U^2+hW-y,\,
                     S(h,U,V,W)-s\bigr)
\]
is isomorphic to $B[T]$, where $T$ is an indeterminate.
\end{lemma}

\begin{proof}
We first make one relation linear, then eliminate $W$ using the unit-ideal
hypothesis, and finally identify the remaining algebra with a polynomial ring
in one variable.
Choose $\alpha,\beta\in B$ with
\begin{equation}
  \alpha h+\beta y=1. \label{eq:bezout}
\end{equation}
Make the polynomial substitution $G=V+UW$, whose inverse is $V=G-UW$.
Modulo the relation $hW=y+U^2$, direct expansion gives
\[
  S(h,U,G-UW,W)=hG^2-2UyG+y^2W.
\]
The first defining relation becomes $hG-yU=x$ modulo the same relation.
Consequently
\begin{equation}
  C\cong B[U,G,W]/\bigl(hG-yU-x,\ hW-y-U^2,\
                       y^2W-s+hG^2-2UyG\bigr). \label{eq:Clinear}
\end{equation}

Put $A=B[U,G]/(hG-yU-x)$ and, in $A$, set
\[
  a=y+U^2,\qquad b=s-hG^2+2UyG.
\]
The last two equations in \eqref{eq:Clinear} are $hW=a$ and $y^2W=b$.
They are compatible, since
\begin{equation}
  y^2a-hb=(hG-yU)^2+y^3-hs=x^2+y^3-hs=0. \label{eq:compat}
\end{equation}
Moreover, if $r=\alpha(1+\beta y)$ and $q=\beta^2$, then
\[
  rh+qy^2=(1-\beta y)(1+\beta y)+\beta^2y^2=1.
\]
It follows that
\begin{equation}
  W_0=ra+qb \label{eq:Wlinear}
\end{equation}
solves both equations: \eqref{eq:compat} gives
$hW_0=(rh+qy^2)a=a$ and $y^2W_0=(rh+qy^2)b=b$.
In $A[W]$ one has
\[
  W-W_0=r(hW-a)+q(y^2W-b),
\]
while $hW-a=h(W-W_0)$ and $y^2W-b=y^2(W-W_0)$. Hence
\[
  (hW-a,\ y^2W-b)=(W-W_0).
\]
Evaluation at $W_0$ therefore gives
$A[W]/(hW-a,y^2W-b)\cong A$. By \eqref{eq:Clinear}, this proves $C\cong A$.

Finally, define a $B$-algebra map $A\to B[T]$ by
\begin{equation}
  U=hT-\beta x,\qquad G=yT+\alpha x. \label{eq:UG}
\end{equation}
Its inverse is given by
$T\mapsto\alpha U+\beta G$.
Indeed, $hG-yU=x$ follows from \eqref{eq:bezout}; in $A$ one has
\[
  h(\alpha U+\beta G)=U+\beta x,\qquad
  y(\alpha U+\beta G)=G-\alpha x,
\]
and substitution of \eqref{eq:UG} gives $\alpha U+\beta G=T$.
These verify both composites. Thus $C\cong A\cong B[T]$.
In the original variables, the parameter is
\[
  T=\alpha U+\beta(V+UW).
\]
Its inverse recovers $U,G$ by \eqref{eq:UG}, then $W$ by
\eqref{eq:Wlinear} and $V=G-UW$.
\end{proof}

\section{The affine-space fibre and the critical point}\label{sec:mainproof}

We first identify the base algebra to which Lemma~\ref{lem:lift} applies.
After obtaining the coordinates on the zero fibre, we use the same polynomial
identities in the ambient ring to locate a critical point on another fibre.

\begin{proof}[Proof of Theorem~\ref{thm:main}]
First, direct expansion of \eqref{eq:xyz} gives
\begin{equation}
  x^2+y^3=hs. \label{eq:cusp}
\end{equation}
For formal $x,y,s$, the translation
\begin{equation}
  X=x+s^3,\qquad Y=y-s^2 \label{eq:shift}
\end{equation}
satisfies
\begin{equation}
  X^2+Y^3=x^2+y^3-sp(x,y,s). \label{eq:shiftidentity}
\end{equation}
Here the terms $s^6$ from the square and cube cancel.

Using independent indeterminates $h,x,y,s$, form the auxiliary algebra
\begin{equation}
  B=\CC[h,x,y,s]/\bigl(h-1-p(x,y,s),\ x^2+y^3-sh\bigr).
  \label{eq:B}
\end{equation}
We use the same symbols for their residue classes in $B$.
This is a new presentation: at this stage $x,y,s$ are independent generators
before passage to the quotient, rather than the polynomials in \eqref{eq:xyz}.
Eliminating $h$ and applying the invertible translation
\eqref{eq:shift} over $\CC[s]$ gives
\begin{equation}
  B\cong\CC[X,Y,s]/(X^2+Y^3-s)\cong\CC[X,Y]. \label{eq:Bplane}
\end{equation}
The inverse presentation is explicitly
\begin{equation}
  s=X^2+Y^3,\quad x=X-s^3,\quad y=Y+s^2,
  \quad h=1+p(x,y,s). \label{eq:Bparam}
\end{equation}

The elements $h,y$ generate the unit ideal in $B$. Indeed, modulo $(h,y)$,
the two defining equations give $x^2=0$ and $1=2s^2x$; squaring the latter
gives $1=0$. Hence $B/(h,y)$ is the zero ring.

Form the algebra in Lemma~\ref{lem:lift} over this $B$, using lower-case
$u,v,w$ for its three added variables. Its presentation is
\[
  C=B[u,v,w]/\bigl(u^3+hv-x,\,-u^2+hw-y,\,
                      S(h,u,v,w)-s\bigr).
\]
In the polynomial ring on $h,x,y,s,u,v,w$, the last three equations eliminate
$x,y,s$. The relation $x^2+y^3-sh$ then becomes the zero polynomial by
\eqref{eq:cusp}, and the other relation in \eqref{eq:B} becomes exactly $F$.
Thus $C\cong R/(F)$. Lemma~\ref{lem:lift} and \eqref{eq:Bplane} now give
\begin{equation}
  R/(F)\cong C\cong B[T]\cong\CC[X,Y,T]. \label{eq:fiber}
\end{equation}

It remains to show that $F$ is not a coordinate. Return to $R$, with
$X=x+s^3$ and $Y=y-s^2$ now interpreted using \eqref{eq:xyz}.
Equations \eqref{eq:cusp}, \eqref{eq:F}, and \eqref{eq:shiftidentity} give
the ambient polynomial identity
\begin{equation}
  X^2+Y^3=s(1+F). \label{eq:ambient}
\end{equation}
At $P=(2,0,-1/2,1/2)$ in the variable order $(h,u,v,w)$, the values are
\[
  x=-1,\quad y=1,\quad s=1,\quad p=2,
  \qquad X=Y=0,\quad F=-1.
\]
Differentiating \eqref{eq:ambient} gives
\[
  2X\,dX+3Y^2\,dY=(1+F)\,ds+s\,dF.
\]
At $P$ this reduces to $0=dF(P)$, so $\grad F(P)=0$.

We now use the ambient critical-point obstruction emphasized by
Shpilrain and Yu~\cite[\S1]{SY}.
If a polynomial automorphism had $F$ as its first component, its Jacobian
matrix at $P$ would have a zero first row. But the chain rule applied to the
automorphism and its polynomial inverse makes that matrix invertible.
This contradiction proves that $F$ is not a coordinate.
\end{proof}

\begin{proof}[Proof of Corollary~\ref{cor:dimension}]
For $n>4$, use the same $F$ in $R[z_1,\ldots,z_{n-4}]$. Equation~\eqref{eq:fiber}
shows that the quotient is a polynomial ring in $n-1$ variables.
All derivatives of $F$ in the added variables vanish, and its original four
derivatives vanish at $P$. Thus the extended polynomial has a critical point
at $(P,0,\ldots,0)$, and the same Jacobian argument rules out a coordinate.
The case $n=4$ is Theorem~\ref{thm:main}.
\end{proof}

\appendix
\section{Explicit coordinates on the zero fibre}\label{app:coordinates}

The proof of Theorem~\ref{thm:main} already supplies both directions of the
isomorphism. We collect them here, including a polynomial certificate for the
unit-ideal condition, so that the parametrization can be used directly.

To make the isomorphism \eqref{eq:fiber} explicit, take in the algebra $B$ of
\eqref{eq:B}
\begin{align}
  \alpha&=1+2s^2x+4s^5,\notag\\
  \beta&=(3s^3-3sy)(1+2s^2x)-4s^4y^2. \label{eq:alphabeta}
\end{align}
The following identity in the free polynomial ring proves
$\alpha h+\beta y=1$ in $B$:
\begin{equation}
  \alpha h+\beta y-1
  =(1+2s^2x)(h-1-p)-4s^4(x^2+y^3-sh). \label{eq:certificate}
\end{equation}

The three polynomial coordinates on $F=0$ are
\begin{equation}
  X=x+s^3,\qquad Y=y-s^2,\qquad
  T=\alpha u+\beta(v+uw), \label{eq:fibercoords}
\end{equation}
with all expressions interpreted by \eqref{eq:xyz} and \eqref{eq:alphabeta}.
Together they define a map from the zero fibre to $\Aff^3_{\CC}$.
Conversely, start from independent $X,Y,T$, reconstruct $s,x,y,h$ by
\eqref{eq:Bparam}, and compute $\alpha,\beta$ by \eqref{eq:alphabeta}. Set
\begin{align}
  u&=hT-\beta x,\qquad g=yT+\alpha x,\notag\\
  w&=\alpha(1+\beta y)(y+u^2)
        +\beta^2(s-hg^2+2uyg),\notag\\
  v&=g-uw. \label{eq:explicitinverse}
\end{align}
These are exactly the inverse maps from Lemma~\ref{lem:lift}.
In particular, they are defined on every point of $\Aff^3_{\CC}$ and every
point of the zero fibre; no open subset has been removed.

\begin{thebibliography}{99}
\small

\bibitem{AM}
S.~S.~Abhyankar and T.~T.~Moh,
\emph{Embeddings of the line in the plane},
J.~Reine Angew. Math. \textbf{276} (1975), 148--166.
\href{https://doi.org/10.1515/crll.1975.276.148}{doi:10.1515/crll.1975.276.148}.

\bibitem{EEO}
M.~El~Kahoui, N.~Essamaoui, and M.~Ouali,
\emph{Interpolation in the automorphism group of a polynomial ring},
Algebra Colloq. \textbf{27} (2020), no.~3, 587--598.
\href{https://doi.org/10.1142/S1005386720000486}{doi:10.1142/S1005386720000486}.

\bibitem{Ghosh}
P.~Ghosh,
\emph{On $\Aff^1$-contractibility of certain simple birational extensions of affine spaces},
\href{https://arxiv.org/abs/2506.09166v4}{arXiv:2506.09166v4} (2026).

\bibitem{GGP}
P.~Ghosh, N.~Gupta, and A.~Pal,
\emph{On embedding of linear hypersurfaces},
\href{https://arxiv.org/abs/2405.07205v3}{arXiv:2405.07205v3} (2024).

\bibitem{GGPthreefold}
P.~Ghosh, N.~Gupta, and A.~Pal,
\emph{On the family of affine threefolds $a(x)y=F(x,z,t)$},
J.~Algebra \textbf{700} (2026), 267--288.
\href{https://doi.org/10.1016/j.jalgebra.2026.04.008}{doi:10.1016/j.jalgebra.2026.04.008};
\href{https://arxiv.org/abs/2403.17397v2}{arXiv:2403.17397v2}.

\bibitem{KVZ}
S.~Kaliman, S.~V\'en\'ereau, and M.~Zaidenberg,
\emph{Simple birational extensions of the polynomial algebra $\CC^{[3]}$},
Trans. Amer. Math. Soc. \textbf{356} (2004), 509--555.
\href{https://doi.org/10.1090/S0002-9947-03-03398-1}{doi:10.1090/S0002-9947-03-03398-1};
\href{https://arxiv.org/abs/math/0104204v2}{arXiv:math/0104204v2}.

\bibitem{Kraft}
H.~Kraft,
\emph{\href{https://numdam.org/item/SB_1994-1995__37__295_0/}{Challenging problems on affine $n$-space}},
S\'eminaire Bourbaki, Vol.~1994/95, Expos\'e~802,
Ast\'erisque \textbf{237} (1996), 295--317.

\bibitem{Maubach}
S.~Maubach,
\emph{The commuting derivations conjecture},
J.~Pure Appl. Algebra \textbf{179} (2003), 159--168.
\href{https://doi.org/10.1016/S0022-4049(02)00198-6}{doi:10.1016/S0022-4049(02)00198-6}.

\bibitem{MaubachMonomial}
S.~Maubach,
\emph{Triangular monomial derivations on $k[X_1,X_2,X_3,X_4]$ have kernel
generated by at most four elements},
J.~Pure Appl. Algebra \textbf{153} (2000), no.~2, 165--170.
\href{https://doi.org/10.1016/S0022-4049(99)00084-5}{doi:10.1016/S0022-4049(99)00084-5}.

\bibitem{Russell}
P.~Russell,
\emph{Simple birational extensions of two dimensional affine rational domains},
Compositio Math. \textbf{33} (1976), no.~2, 197--208.
\href{https://www.numdam.org/item/CM_1976__33_2_197_0/}{Archival full text}.

\bibitem{Sathaye}
A.~Sathaye,
\emph{On linear planes},
Proc. Amer. Math. Soc. \textbf{56} (1976), 1--7.
\href{https://doi.org/10.1090/S0002-9939-1976-0409472-6}{doi:10.1090/S0002-9939-1976-0409472-6}.

\bibitem{SY}
V.~Shpilrain and J.-T.~Yu,
\emph{Embeddings of hypersurfaces in affine spaces},
J.~Algebra \textbf{239} (2001), no.~1, 161--173.
\href{https://doi.org/10.1006/jabr.2000.8677}{doi:10.1006/jabr.2000.8677};
\href{https://arxiv.org/abs/math/0010211v1}{arXiv:math/0010211v1}.

\bibitem{Suzuki}
M.~Suzuki,
\emph{Propri\'et\'es topologiques des polyn\^omes de deux variables complexes,
et automorphismes alg\'ebriques de l'espace $\CC^2$},
J.~Math. Soc. Japan \textbf{26} (1974), 241--257.
\href{https://doi.org/10.2969/jmsj/02620241}{doi:10.2969/jmsj/02620241}.

\end{thebibliography}
\end{document}
